\section{Cartesian Monoidal Categories: A Nice Context for Parameterization}
\label{Sect:CartMon}
%
The goal of this section is to understand nice contexts for parameterization, such as the examples studied in \cref{Sect:Motivation}.
To this end, let $\mathcal{V}^{\mathcal{W}}$ be a braided monoidal category.
If $\mathcal{V}$ is a nice context for parameterization, then it would be reasonable to assume that each object $P \in \mathcal{V}_0$ can be used as the parameter object in at least one way.
This means that there must exists at least one comonoid $M_P = ( P, d_p, e_P )$ in $\mathcal{V}$ from which to obtain a parameterization $\Omega_{(\mathcal{V},\mathcal{W})}( M_P )$.
The comonoid associated with $\mathbb{I}$ should be $( \mathbb{I}, \lambda^{-1}, \mathbb{I} )$, since this is the comonoid used to obtain the underlying $\mathcal{W}$-categories.
It would also be nice if this choice of comonoids respected the monoidal structure on $\mathcal{V}$, in the sense that $M_{P \otimes Q} = M_P \otimes M_Q$.
This makes sense since $\mathcal{V}$ is braided monoidal, from which a monoidal structure can be induced on $\mathbf{Comon}( \mathcal{V} )$.
Now, consider the problem of finding a reparameterization from $\Omega_{(\mathcal{V},\mathcal{W})}( M_Q )$ to $\Omega_{(\mathcal{V},\mathcal{W})}( M_P )$.
In general, this is hard since not every morphism $f \in \mathcal{V}( P, Q )$ lifts to comonoid homomorphism from $M_P$ to $M_Q$.
In a nice context for reparameterization, such a lifting would always exist.
These assumptions can be summarized by requiring the existence of a monoidal section to the forgetful functor from $\mathbf{Comon}( \mathcal{V} )$ to $\mathcal{V}$.

\begin{defi}
    A braided monoidal category $\mathcal{V}$ is a \emph{nice context for parameterization} if there exists a strict monoidal section to the forgetful functor $U : \mathbf{Comon}( \mathcal{V} ) \to \mathcal{V}$.
\end{defi}

\begin{lem}
    \label{Lem:NiceParam}
    A braided monoidal category $\mathcal{V}$ is a nice context for parameterization if and only if the following conditions hold.
    \begin{enumerate}
    \item Each object $P \in \mathcal{V}_0$ admits a comonoid $M_P$.
    \item If $P \in \mathcal{V}_0$ and $Q \in \mathcal{V}_0$, then $M_{P \otimes Q} = M_P \otimes M_Q$.
    \item If $f \in \mathcal{V}( P, Q )$, then $f$ is a comonoid homomorphism from $M_P$ to $M_Q$.
    \item $M_{\mathbb{I}} = ( \mathbb{I}, \lambda^{-1}_{\mathbb{I}}, 1_{\mathbb{I}} )$.
    \end{enumerate}
\end{lem}

\begin{proof}
    First, assume that $\mathcal{V}$ is a nice context for parameterization.
    Then there exists a strict monoidal section $S: \mathcal{V} \to \mathbf{Comon}( \mathcal{V} )$ to the forgetful functor $U: \mathbf{Comon}( \mathcal{V} ) \to \mathcal{V}$.
    For each $P \in \mathcal{V}_0$, define $M_P = S_0( P )$.
    The four properties are satisfied as follows.
    \begin{enumerate}
    \item Let $P \in \mathcal{V}_0$.
          Since $S$ is a section to $U$, then $U_0( M_P ) = U_0( S_0( P ) ) = P$.
          This means that $M_P$ is a comonoid on $P$.
    \item Let $P \in \mathcal{V}_0$ and $Q \in \mathcal{V}_0$.
          Since $S$ is a strict monoidal functor, then by definition of a strict tensorator $M_{P \otimes Q} = S_0( P \otimes Q ) = S_0( P ) \otimes S_0( Q ) = M_P \otimes M_Q$.
    \item Let $f \in \mathcal{V}( P, Q )$.
          Then $S( f )$ is a comonoid homomorphism from $M_P$ to $M_Q$.
          This means that $U( S( f ) )$ lifts to a comonoid homomorphism from $M_P$ to $M_Q$.
          Since $S$ is a section to $U$, then $U( S( f ) ) = f$.
          Then $f$ is a comonoid homomorphism from $M_P$ to $M_Q$.
    \item Since $S$ is strict monoidal, then $M_{\mathbb{I}} = S_0( \mathbb{I} ) = ( \mathbb{I}, \lambda^{-1}_{\mathbb{I}}, 1_{\mathbb{I}} )$.
    \end{enumerate}
    This concludes the first direction of the proof.
    Next, assume that $\mathcal{V}$ is a braided monoidal category satisfying properties (1) through to (4).
    By property (1), there exists an assignment $S_0: \mathcal{V}_0 \to \mathbf{Comon}( \mathcal{V} )_0$ such that $S_0: P \mapsto M_P$.
    By property (3), this extends to an assignment on morphisms, such that $S_{P,Q}: f \mapsto f$.
    Let $S: \mathcal{V} \to \mathbf{Comon}( \mathcal{V} )$ denote the full assignment from $\mathcal{V}$ to $\mathbf{Comon}( \mathcal{V} )$.
    This is clearly a functor since function composition is the same in $\mathbf{Comon}( \mathcal{V} )$ as in $\mathcal{V}$.
    It can then be shown that $S$ is a section to $U: \mathbf{Comon}( \mathcal{V} ) \to \mathcal{V}$.
    \begin{itemize}
    \item \textbf{Objects}.
          Let $P \in \mathcal{V}_0$.
          Since $M_P$ is a comonoid on $P$, then $U_0( S_0( P ) ) = U_0( M_P ) = P$.
    \item \textbf{Morphisms}.
          If $f \in \mathcal{V}( P, Q )$, then $U( S( f ) ) = U( f ) = f$.
    \end{itemize}
    Then $U \circ S = 1_{\mathcal{V}}$.
    It remains to be shown that $S$ is a strict monoidal functor.
    \begin{itemize}
    \item \textbf{Preserves Products}.
          If $f \in \mathcal{V}( P, Q )$ and $g \in \mathcal{V}( R, S )$, then $S( f \otimes g ) = f \otimes g = S( f ) \otimes S( g )$.
    \item \textbf{Strict Coherence}.
          We first show that $S$ preserves the left unitor.
          If $\lambda$ is the left unitor for $\mathcal{V}$, then $\lambda$ is also the left unitor for $\mathbf{Comon}( \mathcal{V} )$ and $S( \lambda_P ) = \lambda_P$ for each $P \in \mathcal{V}_0$.
          A similar argument holds for the right unitor of $\mathcal{V}$ and the associator of $\mathcal{V}$.
    \item \textbf{Strict Tensorator}.
          By property (2), $S( P \otimes Q ) = M_{P \otimes Q} = M_P \otimes M_Q = S( P ) \otimes S( Q )$ for each $P \in \mathcal{V}_0$ and $Q \in \mathcal{V}_0$.
    \item \textbf{Strict Unitor}.
          By property (4), $S( \mathbb{I} ) = M_{\mathbb{I}} = ( \mathbb{I}, \lambda^{-1}_{\mathbb{I}}, 1_{\mathbb{I}} )$.
    \end{itemize}
    Then $S$ is a strict monoidal section to the forgetful functor.
\end{proof}

It turns out that a symmetric monoidal category $\mathcal{V}$ is a nice context for parameterization if and only if $\mathcal{V}$ is Cartesian monoidal.
This can be seen as a variation on Fox's theorem.
To prove the first direction, it suffices to prove that the strictness of $S: \mathcal{V} \to \mathbf{Comon}( \mathcal{V} )$ ensures that $\mathbb{I}$ is terminal, and that the comonoids chosen by $S$ can be used to define uniform copying and deleting.
To prove the second direction, it suffices to show that the uniform copying and deleting exhibited by Cartesian monodial categories pick out comonoids satisfying the properties in \cref{Lem:NiceParam}.

\begin{thm}
    \label{Thm:NiceParam}
    A symmetric monoidal category $\mathcal{V}$ is a nice context for parameterization if and only if $\mathcal{V}$ is Cartesian monoidal category.
\end{thm}

\begin{proof}
    First, assume that $\mathcal{V}$ is a nice context for parameterization.
    Then there exist a strict monoidal section $S: \mathcal{V} \to \mathbf{Comon}( \mathcal{V} )$ to the forgetful functor $U: \mathbf{Comon}( \mathcal{V} ) \to \mathcal{V}$.
    It must first be shown that $\mathbb{I}$ is terminal.
    To this end, let $P \in \mathcal{V}$ and $f \in \mathcal{V}( P, \mathbb{I} )$.
    Since $S$ is strict monoidal, then $S_0( \mathbb{I} ) = ( \mathbb{I}, \lambda_{\mathbb{I}}^{-1}, 1_{\mathbb{I}} )$.
    Then $f = U( S( f ) )$ is a comonoid homomorphism from $S_0( P ) = ( P, d, e )$ to $S_0( \mathbb{I} ) = ( \mathbb{I}, \lambda_{\mathbb{I}}^{-1}, 1_{\mathbb{I}} )$.
    Then by definition of a comonoid homomorphism, $e = 1_{\mathbb{I}} \circ f = f$.
    Since $f$ was arbitrary, then $\mathcal{V}( P, \mathbb{I} ) = \{ e \}$.
    Since $P$ was arbitrary, then $\mathbb{I}$ is terminal.
    Next, it must be shown that $\mathcal{V}$ has uniform copying and uniform deleting.
    \begin{itemize}
    \item \textbf{Uniform Copying}.
          For each $P \in \mathcal{V}_0$, let $\Delta_P$ denote the comultiplication in $S_0( P )$.
          Then $\Delta: \mathcal{V} \Rightarrow \mathcal{V} \otimes \mathcal{V}$ defines an unnatural transformation.
          To show that $\Delta$ is natural, let $f \in \mathcal{V}( P, Q )$ for some $P \in \mathcal{V}_0$ and $Q \in \mathcal{V}_0$.
          Then $S( f )$ is a comonoid homomorphism.
          In particular, $\Delta_Q \circ f = ( f \otimes f ) \circ \Delta_P$.
          Since $f$ was arbitrary, then $\Delta$ is natural.
          Since $S$ is a strict monoidal functor, then $\Delta_\mathbb{I} = \lambda_{\mathbb{I}}^{-1}$.
          It remains to be shown that $\Delta$ defines uniform copying on $\mathcal{V}$.
          Let $P \in \mathcal{V}_0$ and $Q \in \mathcal{V}_0$.
          Since $S$ is a strict monoidal functor, then $S( P \otimes Q ) = S( P ) \otimes S( Q )$.
          Since $\Delta_{P \otimes Q}$ is the comultiplication for $S( P ) \otimes S( Q )$, then \cref{Eqn:MonoidProd} holds in $\mathcal{V}$.
          Since $P$ and $Q$ were arbitrary, then $\Delta$ defines uniform copying on $\mathcal{V}$.
    \item \textbf{Uniform Deleting}.
          Since $\mathbb{I}$ is terminal, then there exists a unique natural transformation $\epsilon: \mathcal{V} \Rightarrow \mathbb{I}$.
          Moreover, if $P \in \mathcal{V}_0$ and $( P, d, e ) = S_0( P )$, then $\epsilon_P = e$ by uniqueness.
    \end{itemize}
    By construction, $( P, \Delta_P, \epsilon_P ) = S_0( P )$ is a comonoid for each $P \in \mathcal{V}_0$.
    Then $\mathcal{V}$ is Cartesian monoidal by \cref{Prop:Cartesian}.
    This concludes the first direction of the proof.
    Next, assume that $\mathcal{V}$ is Cartesian monoidal.
    Then $\mathcal{V}$ admits uniform copying $\Delta: \mathcal{V} \Rightarrow \mathcal{V} \otimes \mathcal{V}$ and uniform deleting $\epsilon: \mathcal{V} \Rightarrow \mathbb{I}$.
    Since $\Delta$ and $\epsilon$ are compatible, then each object $P \in \mathcal{V}_0$ admits a comonoid $( P, \Delta_P, \epsilon_P )$.
    These comonoids satisfy the properties of \cref{Lem:NiceParam}.
    \begin{enumerate}
    \item For each $P \in \mathcal{V}_0$, define $M_P = ( P, \Delta_P, \epsilon_P )$.
    \item Let $P \in \mathcal{V}_0$ and $Q \in \mathcal{V}_0$.
          Since $\mathcal{V}$ is Cartesian monoidal, then $\mathbb{I}$ is terminal.
          It follows by uniqueness that $\epsilon_{P \otimes Q} = \lambda_{\mathbb{I}}^{-1} \circ ( \epsilon_P \otimes \epsilon_Q )$.
          Moreover, since $\Delta$ defines uniform copying, then \cref{Eqn:MonoidProd} holds in $\mathcal{V}$.
          Then by definition, $M_{P \otimes Q} = M_P \otimes M_Q$.
    \item Let $f \in \mathcal{V}( P, Q )$.
          Since $\Delta$ is a natural transformation, then $( f \otimes f ) \circ \Delta_P = \Delta_Q \circ f$.
          Since $\mathcal{V}$ is Cartesian monoidal, then $\mathbb{I}$ is terminal.
          Then $\epsilon_P = \epsilon_Q \circ f$ by uniqueness.
          Then $f$ is a comonoid homomorphism from $M_P$ to $M_Q$.
    \item By definition, $\Delta_{\mathbb{I}} = \lambda_{\mathbb{I}}^{-1}$ and $\epsilon_{\mathbb{I}} = 1_\mathbb{I}$.
          Then $M_{\mathbb{I}} = ( \mathbb{I}, \lambda_{\mathbb{I}}^{-1}, 1_{\mathbb{I}} )$.
    \end{enumerate}
    Then $\mathcal{V}$ is a nice context for parameterization by \cref{Lem:NiceParam}.
\end{proof}

It turns out that even if $\mathcal{V}$ is a nice context for parameterization, it may not be the case that evaluation is possilbe.
It was shown in~\cref{Ex:NoParams} that every meet-semilattice with at least two elements defines a Cartesian monoidal category in which parameter objects never have elements.
In contrast, topological spaces define a category in which parameters always have elements.
Similar, the category of sub-lattices is also Cartesian monoidal and admits evaluation, giving rise to interesting examples from classical program analysis.
It should be noted that this is not the standard monoidal closed structure on \textbf{SupLat}, since \textbf{SupLat} is not closed as a Cartesian monoidal category.

\begin{exa}[Parameterized Rotations]
    \label{Ex:CartRot}
    Since $\mathbf{Top}$ is a Cartesian monoidal category, then by \cref{Thm:NiceParam} the category $\mathbf{Top}$ is a nice context for parameterization.
    This means that the categories described in the motivating example inherit many nice properties from $\mathbf{Top}$.
    To this end, let $\mathcal{V} = \mathbf{Top}$, $\mathcal{W} = \mathbf{Set}$, and $\mathcal{C}^{\mathcal{V}} = \mathbb{C}\mathbf{FVect}^{\mathcal{V}}$.
    Then for each $P \in \mathbf{Top}$, the category $\mathcal{C}_P := \Omega_{(\mathcal{V},\mathcal{W})}( M_P )( \mathcal{C}^{\mathcal{V}} )$ is the ordinary category whose objects are vector spaces with morphisms from $V$ to $W$ given by the continuous maps from $P$ into $\mathcal{C}^{\mathcal{V}}( V, W )$.
    Since $\mathcal{V}$ is a nice context for parameterization, then all moprhisms in $\mathcal{V}$ induce reparameterizations.
    In particular, the structural maps which make $\mathcal{V}$ Cartesian will also induce reparameterizations.
    \begin{itemize}
    \item \textbf{Symmetries}.
          Let $P = X \times Y$, $Q = Y \times X$, and $\sigma_{X,Y}$ denote the symmetry from $P$ to $Q$ which sends each tuple $( x, y )$ to $( y, x )$.
          Then $F = \Omega_{(\mathcal{V},\mathcal{W})}( \sigma_{X,Y} )( \mathcal{C} ): \mathcal{C}_Q \to \mathcal{C}_P$ is the functor which maps each morphsim $f \in \mathcal{C}_Q( V, W )$ to the morphism $( x, y ) \mapsto f( y, x )$ in $\mathcal{C}_P( V, W )$.
          This provides a framework to discuss whether a morphism is symmetric in its parameters.
          For example, if $X = Y$, then $f( x, y ) = f( y, x )$ if and only if $f = F( f )$.
    \item \textbf{Projections}.
          Let $P = X \times Y$.
          Then there exists a projection map $\pi_X: P \to X$ such that $\pi_X( x, y ) = x$.
          Then $F = \Omega_{(\mathcal{V},\mathcal{W})}( \pi_X )( \mathcal{C} ): \mathcal{C}_X \to \mathcal{C}_P$ is the functor which maps each morphism $f \in \mathcal{C}_X( V, W )$ to the morphism $( x, y ) \to f( x )$ in $\mathcal{C}_P( V, W )$.
          In other words, $F$ picks out the subcategory of $\mathcal{C}_P$ consisting of the parameterized morphisms in $\mathcal{C}_P$ which do not depend on $Y$.
    \item \textbf{Copying}.
          Let $X \in \mathcal{V}_0$ and $P = X \times X$.
          Then there exists a copying map $\Delta_X: X \to P$ such that $\Delta_X( x ) = ( x, x )$.
          Then $\Omega_{(\mathcal{V},\mathcal{W})}(\Delta_X )( \mathcal{C} ): \mathcal{C}_P \to \mathcal{C}_X$ is the functor which maps each morphsim $f \in \mathcal{C}_P( V, W )$ to the morphism $x \mapsto f( x, x )$ in $\mathcal{C}_X( V, W )$.
          In the case of parameterized quantum circuits, this provides a framework to interpret circuits in which a rotation gate depends on the same parameter more than once.
    \end{itemize}
    It should be noted that reparameterizations already appear in the existing work on parameterized quantum circuit analysis, though these results were not stated in the language of enriched category theory.
    For example, the authors of~\cite{RossWesley2025} consider circuits parameterized by $P = \mathbb{R}^k$.
    These circuits are then precomposed with a linear automorphism $f: P \to P$ to reduce parameterized equivalence checking with rational coefficients to parameterized equivalence checking with integer coefficients.
    Since $f$ is a linear map between finite-dimensional vector space, then clearly $f$ is continuous.
    This means that the reparameterization induced by $f$ is indeed the isomorphism $\Omega_{(\mathcal{V},\mathcal{W})}( f )( \mathcal{C} ): \mathcal{C}_P \to \mathcal{C}_P$.
\end{exa}

\begin{exa}[Discrete-Time Processes]
    Let $\mathcal{P}: \mathbf{Set} \to \mathbf{SupLat}$ denote the covariant power-set functor.
    It was shown in~\cite{Little2025}, that non-determinism in program semantics could be understood in terms of \textbf{SupLat}-enrichment.
    For example, assume that the semantics of a deterministic programming language are given by some locally small category $\mathcal{C}$.
    This means that $\mathcal{C}( X, Y )$ is the set of deterministic programs which convert inputs of type $X$ into outputs of type $Y$.
    Given two deterministic programs $f \in \mathcal{C}( X, Y )$ and $g \in \mathcal{C}( X, Y )$, a non-deterministic program is allowed to choose between $f$ and $g$, independent of the input data~\cite{10.1145/321420.321422}.
    Non-determinism is often used as a programming abstraction to model interactions with an external environment, such as requesting an input from a user.
    As described in~\cite{Little2025}, non-determinism can be introduced to $\mathcal{C}$ via change-of-base through the power-set functor.
    The interpretation of $\mathcal{D} = \mathcal{P}_*( \mathcal{C} )$ is as follows.
    \begin{itemize}
    \item \textbf{Objects}.
          Since $\mathcal{D}_0 = \mathcal{C}_0$, then the collection of types remains unchanged.
    \item \textbf{Morphisms}.
          If $\widetilde{f} \in \mathcal{D}( X, Y )$, then $\widetilde{f} \subseteq \mathcal{C}( X, Y )$.
          This means that $\widetilde{f}$ is a set of deterministic programs.
          Each $f \in \widetilde{f}$ is a possible outcome when $\widetilde{f}$ is executed.
          For example, $\widetilde{f} = \{ f, g \}$ is the non-deterministic program which chooses between executing either $f$ or $g$.
    \item \textbf{Determinism}.
          There is an inclusion $j: \mathcal{C}( X, Y ) \hookrightarrow \mathcal{D}( X, Y )$ such that $j: f \mapsto \{ f \}$.
          This inclusion maps each program in $\mathcal{C}( X, Y )$ to a singleton set, and can be seen as picking out the deterministic programs in $\mathcal{D}( X, Y )$.
    \item \textbf{Composition}.
          If $\widetilde{f} \in \mathcal{D}( X, Y )$ and $\widetilde{g} \in \mathcal{D}( Y, Z )$, then $\widetilde{g} \circ \widetilde{f} = \{ g \circ f : f \in \widetilde{f} \land g \in \widetilde{g} \}$.
          Then the composition of $\widetilde{f}$ with $\widetilde{g}$ is composition of all possible outcomes from $\widetilde{f}$ with all possible outcomes from $\widetilde{g}$.
    \end{itemize}
    An abstraction of discrete-time systems can then be obtained by parameterizing $\mathcal{D}$ by the extended natural numbers $L = \mathbb{N} \cup \{ \infty \}$.
    To see how this works, let $F = \Omega_{(\mathbf{SupLat},\mathbf{Set})}( L )$ and $\mathcal{D}_L = F( \mathcal{D} )$.
    Then each $\widetilde{f} \in \mathcal{D}_L( X, Y )$ is a monotonic function from $L$ to $\mathcal{D}( X, Y )$.
    Each $f \in \widetilde{f}( n )$ is a possible outcome from executing $\widetilde{f}$ at some time between $0$ and $n$.
    This can be used to model behaviours such as error recovery, as illustrated in the following toy model.
    In this example, assume that $a \in \mathcal{D}( X, Y )$ is the desired behaviour of a system, and $b \in \mathcal{D}( X, y )$ is a possible error.
    Further assume that $c \in \mathcal{D}( Y, Y )$ is an error-recovery protocol, such that $c \circ a = c \circ b = a$.
    We will model the system by some $\widetilde{f} \in \mathcal{D}_L( X, Y )$ such that $\widetilde{f}( n ) \subseteq \{ a, b \}$ for all $n \in L$, $\widetilde{f}( t ) = \{ a, b \}$ for some $t \in \mathbb{N}$.
    This means that $\widetilde{f}$ will fail after some unknown time $t$.
    If we assume that assume that $c$ is fault-tolerant (i.e., it will never fail), then we may lift $c$ to $\widetilde{g} = j( c ) \circ \epsilon_L$.
    That is, $\widetilde{g}( n ) = \{ c \}$ for all $n \in L$.
    Since $( \widetilde{g} \star \widetilde{f} )( n ) = \widetilde{g}( n ) \circ \widetilde{f}( n )$, then it can be proven that $( \widetilde{g} \star \widetilde{f} )( n ) = j( a ) = \{ a \}$ for all $n \in L$.
    In other words, the process $\widetilde{g} \star \widetilde{f}$ can recover from all of the faults encountered by $\widetilde{f}$.
\end{exa}
